\documentclass[12pt]{article}
\usepackage[utf8]{inputenc}
\usepackage[T1]{fontenc}
\usepackage{hyperref}
\title{Mo Language Resources for K-12 Curriculum Development}
\author{Compiled by Claude Sonnet 5.5 (claude-sonnet-5-5) for Andr\'as Kornai}
\date{}
\begin{document}
\maketitle

\section{Introduction}
Mo (ISO 639-3 \texttt{wkd}, Glottocode wakd1237) is also called Wakde. It is an Austronesian language of the Sarmi--Jayapura Bay group, spoken on the coast and on Wakde Island in Papua, Indonesia. Wikipedia reports about 550 native speakers (2005). Do not confuse it with Mor (\texttt{mhz}), a separate language near Nabire. The search found very little material beyond encyclopedic and mission-statistics pages.

\section{Primary Educational Resources}
No primers, textbooks or school materials in Mo turned up in search results.

\section{Linguistic and Reference Materials}
\begin{itemize}
\item Wikipedia, Wakde language: \url{https://en.wikipedia.org/wiki/Wakde_language}
\item Wikipedia disambiguation page ``Mo language'': \url{https://en.wikipedia.org/wiki/Mo_language}
\end{itemize}
No dictionary or grammar was found.

\section{Religious and Community Resources}
Joshua Project has a profile of the Wakde people. It gives a population of about 2,000 and lists two dialects, Ma Wihi and Arti. It also labels the language ``Papuan'', which conflicts with the Austronesian classification used elsewhere. \url{https://joshuaproject.net/people_groups/15831}

\section{Digital and Audio Resources}
No audio, video, corpus or software resources were found.

\section{Curriculum Development Context}
The language is spoken by adults only and is not known to be taught in schools. A curriculum would have to begin with new documentation (word lists, recordings and an orthography) made together with the community. No existing educational base was found.
\end{document}